PP has appeared in redistricting litigation and statutory frameworks in multiple
contexts.

\subsection{Court Usage}

\paragraph{North Carolina legislative redistricting (2017--2020).}
Expert reports submitted in North Carolina legislative redistricting litigation
(\textit{Common Cause v.\ Rucho}, later renamed after the state defendants)
cited PP values for proposed and enacted plans.
\citet{chen2017expert} reports PP values for North Carolina Senate districts
in the range $0.10$--$0.35$, consistent with our bisect estimates for NC
congressional districts.

\paragraph{Gill v.\ Whitford (2018).}
The Wisconsin gerrymandering case \textit{Gill v.\ Whitford}, 585 U.S.\ 48 (2018),
produced a substantial record of expert reports.
Compactness metrics including PP were submitted to document that the challenged
legislative plan was less compact than algorithmic alternatives~\citep{gill2018}.

\paragraph{Rucho v.\ Common Cause (2019).}
\textit{Rucho v.\ Common Cause}, 588 U.S.\ 684 (2019), held that partisan
gerrymandering claims are non-justiciable under federal law but explicitly
confirmed that state law compactness standards---including PP---remain
judicially enforceable~\citep{rucho2019}.

\subsection{Statutory Requirements}

\paragraph{Ohio H.B.\ 1 (2015).}
Ohio's 2015 redistricting reform required the General Assembly to maximise
compactness using a PP-based criterion.
The Ohio Redistricting Commission has applied PP in its evaluation framework
for the 2020 redistricting cycle.

\paragraph{Iowa statutory compactness.}
Iowa Code \S 42.4 requires districts to be ``as compact as possible,''
which the Legislative Services Agency has operationalised using perimeter-to-area
ratios equivalent to a PP calculation.

\paragraph{Princeton Gerrymandering Project.}
The Princeton Gerrymandering Project grades redistricting plans on a letter
scale that includes PP as a component of its ``shape'' score~\citep{pgp2022}.
Plans are graded A--F; bisect ratio-optimal plans would score in the B--A range
based on Table~\ref{tab:pp_main}.

\subsection{Limitations for Court Filings}

PP must be accompanied by the projection specification when submitted as
court evidence. An expert report stating ``mean PP $= 0.22$'' without specifying
EPSG:5070 is open to challenge from opposing experts who may compute PP in
a different coordinate system and obtain different values.
We recommend including the following disclosure in all court filings:
\begin{quote}
  \emph{All Polsby-Popper scores were computed using district polygon geometries
  projected to the Albers Equal Area Conic coordinate system (EPSG:5070) using
  the TIGER/Line 2020 Census tract boundaries released by the U.S.\ Census Bureau.
  Perimeter measurements use Euclidean distance between projected polygon vertices.
  No boundary smoothing was applied.}
\end{quote}
